\chapter{Nocturia}

\section*{Definition}
Nocturia is the complaint that an individual must wake at night one or more times to void, with two or more episodes per night considered clinically significant. It is a common cause of sleep disruption and diminished quality of life.

\section*{Pathophysiology}
Nocturia can result from (1) nocturnal polyuria (excessive urine production at night), (2) reduced bladder capacity (structural or functional), (3) global polyuria (24-hour overproduction), or (4) mixed causes. Nocturnal polyuria involves an altered circadian rhythm of antidiuretic hormone (ADH) secretion or responsiveness, leading to excess nighttime urine.

\section*{Causes}
\begin{itemize}
  \item \textbf{Urologic:} Benign prostatic hyperplasia (BPH), overactive bladder, detrusor overactivity, bladder outlet obstruction, urinary tract infection, bladder calculi, interstitial cystitis, prostate cancer, urethral stricture
  \item \textbf{Endocrine/metabolic:} Diabetes mellitus (osmotic diuresis), diabetes insipidus (ADH deficiency or resistance), primary polydipsia, hypercalcemia
  \item \textbf{Cardiovascular:} Congestive heart failure (mobilization of peripheral edema recumbent), venous insufficiency
  \item \textbf{Sleep disorders:} Obstructive sleep apnea (ANP/BNP release from atrial stretch), insomnia, restless legs syndrome
  \item \textbf{Renal:} Chronic kidney disease, nephrotic syndrome (diminished concentrating ability)
  \item \textbf{Medication-induced:} Diuretics (especially evening dosing), caffeine, alcohol, lithium, SSRI, calcium channel blockers
  \item \textbf{Neurological:} Parkinson disease, multiple sclerosis, spinal cord injury (neurogenic bladder)
  \item \textbf{Aging:} Decreased bladder capacity, decreased ADH secretion, altered sleep architecture
\end{itemize}

\section*{When to See a Doctor}
See your doctor if:
\begin{itemize}
  \item Nocturia significantly impairs sleep quality or daytime functioning
  \item There is dysuria, hematuria, or acute onset of diuresis with polydipsia
  \item Associated with orthopnea, leg edema, or paroxysmal nocturnal dyspnea (suspect heart failure)
\end{itemize}

\section*{Related Symptoms}
\begin{itemize}
  \item Urinary frequency, urgency, hesitancy, weak stream, dribbling, dysuria, polyuria, polydipsia, insomnia, daytime fatigue
\end{itemize}
\textbf{Important:} A voiding diary (timed volumes for 72 hours) is the most useful diagnostic tool for distinguishing nocturia subtypes.

\section*{Self-Care / Home Management}
\begin{itemize}
  \item Reduce evening fluid intake, especially 2--3 hours before bed
  \item Avoid caffeine, alcohol, and diuretic beverages after dinner
  \item Elevate legs in the evening to reduce fluid redistribution
  \item Double-void before bed to empty the bladder completely
\end{itemize}

\section*{How Your Doctor Will Evaluate You}
\begin{itemize}
  \item Voiding diary (frequency, volumes, timing) for 3 days
  \item Physical exam: digital rectal exam (prostate), pelvic exam, lower extremity edema, abdomen (distended bladder)
  \item Urinalysis to exclude infection, hematuria, glucosuria
  \item Serum chemistry: glucose, HbA1c, calcium, creatinine, BNP if heart failure suspected
  \item Post-void residual (bladder scan) to assess for retention
  \item Uroflowmetry if bladder outlet obstruction suspected
\end{itemize}

\section*{Treatment Options}
\begin{itemize}
  \item \textbf{Nocturnal polyuria:} Desmopressin (DDAVP) -- synthetic ADH analog; monitor for hyponatremia
  \item \textbf{BPH/bladder outlet obstruction:} Alpha-blockers (tamsulosin), 5-alpha-reductase inhibitors (finasteride)
  \item \textbf{Overactive bladder:} Anticholinergics (oxybutynin, tolterodine) or beta-3 agonists (mirabegron)
  \item \textbf{Heart failure:} Optimize diuretic timing (morning dosing) and heart failure management
  \item \textbf{Diabetes:} Improve glycemic control; treat diabetes insipidus with desmopressin
  \item \textbf{Sleep apnea:} CPAP therapy frequently resolves nocturia
\end{itemize}

\clearpage
